\begin{exercisechunk}
\item \textbf{The angle between tasks controls the eigengap.}
\label[exercise]{ex:l13-moment-eigenspace}
\exerciseprofile{2}{2}{2}{\excalculation\quad\exinterpretation}{%
Relate geometric diversity among two task heads to stability of population
moment recovery.}
Let $k=2$, and sample two source tasks uniformly. Their unit-norm heads are
\[
 a_1=\begin{pmatrix}1\\0\end{pmatrix},
 \qquad
 a_2=\begin{pmatrix}\cos\alpha\\\sin\alpha\end{pmatrix},
 \qquad 0<\alpha\le\frac\pi2.
\]
\begin{enumerate}[(a),beginpenalty=10000]
\item Compute
$D=(a_1a_1^\top+a_2a_2^\top)/2$ and show that its eigenvalues are
\[
 \lambda_{\max}(D)=\frac{1+\cos\alpha}{2},
 \qquad
 \lambda_{\min}(D)=\frac{1-\cos\alpha}{2}.
\]
\item Under the Gaussian model of
\cref{prop:gaussian-task-moment}, show that the gap between the second and
third eigenvalues of $M_\star$ is $1-\cos\alpha$.
\item Determine the leading-order behavior of this eigengap as
$\alpha\downarrow0$. Explain why collecting more observations from two nearly
parallel tasks is statistically different from observing two well-separated
tasks.
\end{enumerate}
\begin{hints}
The trace and determinant of $D$ determine its two eigenvalues. Use
$1-\cos\alpha\sim\alpha^2/2$ for the last part.
\end{hints}
\end{exercisechunk}
